\subsection*{Learning to point and spell in conversational reading}
We used CoQA (Reddy, Chen and Manning, 2019, \url{https://doi.org/10.1162/tacl_a_00266}). We retained complete printable ASCII stories up to 320 words, a current question up to 32 words and up to the last 48 words of the previous human question and answer. Passage hashes determined the reserved training-story split. We excluded 372 official development turns on passages duplicated in the fitting set. The resulting sets comprised 73,714 training turns from 4,476 stories, 4,413 reserved turns from 268 stories and 5,588 development turns from 355 stories. All inputs used a single print face with known grid geometry and visible whitespace. Gold passage evidence was available only during training. At evaluation each turn received the original preceding human exchange. This measures reading a supplied conversational context.

A recurrent state $h_t$ attends to the pictured inputs and receives its own previous word image during generation. The decision distribution has three outcomes: end the reply, point to a word, or spell it. Pointer candidates comprise visible words and their nonempty one-cell boundary trims. For a reference word image $w$ with matching physical candidates $M(w)$, the pointer loss is
\[
 -\log \sum_{j\in M(w)}p(j\mid h_t,\mathrm{point}).
\]
A missing word trains a glyph-sequence likelihood including its end symbol. The speller selects from 89 distinct printed training shapes. The route likelihood includes the answer stop. The source rationale provides an auxiliary endpoint loss with weight 0.20. Lean checks expressibility for reachable word sequences and additivity of matching pointer mass. These statements do not bound the probability learned for a valid sequence.

For the second writer we standardized 115-dimensional measured word-image features and normalized them to unit length to obtain $u(w)$. The predicted vector $v(h_t)$ was trained using cosine distance and a contrastive loss at temperature 0.08 that sums probability over identical word images in a minibatch. The spelling decoder first learned to reconstruct images of 122,940 distinct eligible training word surfaces over 5,000 updates of 192 samples. Both writers then saw the same 8,000 batches of 32 conversation turns. Training used AdamW at 0.00075, gold preceding answer words, a maximum answer length of 16 words and up to 24 cells per word. Held targets exceeding the limits remained in the denominator. Decoding was greedy. Only stopped outputs could receive credit.

A single RTX 4090 D took 132.6 seconds and 15.0 kJ for the direct dialogue fit. The second fit took 192.1 seconds and 20.5 kJ, in addition to 20.5 seconds and 2.20 kJ for word reconstruction. Peak measured tensor allocation during these fits was under 0.44 GB. Each model inherits a frozen 1,900,594-parameter encoder and the passage network initialization. The deployed totals are 2,447,719 and 2,510,170 parameters. Device readings exclude earlier training, data preparation, evaluation and host power. A single seed and development-informed treatment limit inference about architecture and compute requirements.

For the registered wording test, the multiword reference was absent as a case-insensitive substring of each input field and had never been a complete normalized training answer. A stricter recount required at least two words after normalization and an absent normalized phrase. We also joined the four original human references to each eligible development turn by passage hash, turn, question, history and primary answer, and used CoQA's official leave-one-out scoring. The date substitution was a separate paired test on unique years absent from the question and history.
